\documentclass[12pt]{jlreq}
\usepackage{amsmath, amssymb}

\newcommand{\C}{\mathbb{C}}
\newcommand{\R}{\mathbb{R}}
\newcommand{\Z}{\mathbb{Z}}
\newcommand{\Tr}{\operatorname{Tr}}
\newcommand{\Hom}{\operatorname{Hom}}
\newcommand{\GL}{\operatorname{GL}}
\newcommand{\rank}{\operatorname{rank}}
\newcommand{\Id}{\operatorname{Id}}
\newcommand{\Ker}{\operatorname{Ker}}

\begin{document}

束 $(L, \le, \wedge, \vee)$ に関する次の性質 (C) を考える：
\begin{enumerate}
  \item[(C)] 任意の $a, b, c, d, x \in L$ に対して，$a \le b \vee x$ かつ $x \wedge c \le d$ ならば，$a \wedge c \le b \vee d$.
\end{enumerate}
次の問いに答えよ。

\begin{enumerate}
  \item $L$ が分配束ならば，束 $L$ は性質 (C) をみたすことを示せ。
  \item $L, M$ を束とし，写像 $f$ を $L$ から $M$ の上への束の準同型とする。このとき，$L$ が性質 (C) をみたすならば，$M$ も性質 (C) をみたすことを示せ。
\end{enumerate}
また，次の性質 (C') を考える：
\begin{enumerate}
  \item[(C')] 任意の $a, b, c, d, x \in L$ に対して，$a \le b \vee x$ かつ $x \wedge c \le d$ かつ $b \le c$ ならば，$a \wedge c \le b \vee d$.
\end{enumerate}

\begin{enumerate}
  \setcounter{enumi}{2}
  \item $L, M$ を束とし，写像 $f$ を $L$ から $M$ の上への束の準同型とする。$L$ が性質 (C') をみたすならば，$M$ も性質 (C') をみたすことを示せ。
\end{enumerate}

\end{document}